\begin{table}[htbp]
\centering
\footnotesize
\caption{Co-citation clusters on LLM research}
\label{tab:E}
\begin{tabularx}{\textwidth}{|l|Y|c|Y|}
\hline
\textbf{Cluster No and colour} & \textbf{Cluster label} & \textbf{Number of articles} & \textbf{Representative publication} \\ \hline
1 (red) & Edge LLM surveys, collaborative serving, and NLP benchmarks & 128 & Zhang (2024), EdgeShard: Efficient LLM Inference via Collaborative Edge Computing \\ \hline
2 (green) & Quantisation, pruning, and serving efficiency, with foundation models & 95 & Dettmers (2023), QLoRA: Efficient Finetuning of Quantized LLMs \\ \hline
3 (blue) & Mixed background: efficient inference, federated and parameter-efficient tuning, early applications & 58 & Alizadeh (2024), LLM in a flash: Efficient Large Language Model Inference with Limited Memory \\ \hline
4 (yellow) & General background: reasoning, edge computing visions, LLM agents, compression basics & 57 & Wei (2022), Chain-Of-Thought Prompting Elicits Reasoning in Large Language Models \\ \hline
5 (purple) & Generative AI in supply chain and operations management & 6 & Jackson (2024), Generative artificial intelligence in supply chain and operations management: a... \\ \hline
\end{tabularx}
\par\vspace{2pt}{\normalsize Source: Author interpretation based on VOSviewer analysis.}
\end{table}
